% TOP_PROBLEM: {"id": 132, "title": "Collinear 4-tuples Force Collinear 5-tuples", "category_id": 6, "category": {"id": 6, "name": "geometry", "display_name": "Geometry", "description": "Euclidean and non-Euclidean geometry, geometric structures.", "slug": "geometry", "order_index": 6, "created_at": "2026-07-31T15:26:25.671Z"}, "status": "open"}
% FIRST_POSTED: null
% ENTRY_KIND: statement_only
% Imported from: batch03.tex; UnsolvedMath ID 132
\documentclass[11pt]{article}
\usepackage[margin=25mm]{geometry}
\usepackage{fontspec}
\setmainfont{Latin Modern Roman}
\usepackage{amsmath,amssymb,mathtools}
\usepackage{unicode-math}
\setmathfont{Latin Modern Math}
\usepackage{xurl,hyperref}
\hypersetup{colorlinks=true,urlcolor=blue}
\setcounter{secnumdepth}{0}
\setlength{\parindent}{0pt}
\setlength{\parskip}{5pt}
\setlength{\emergencystretch}{3em}
\newcommand{\sref}[2]{\hyperlink{source:#1:#2}{[S#2]}}
\newcommand{\cataloguescope}[1]{\paragraph{Recorded target.}#1}
\begin{document}
% UnsolvedMath ID 132; source catalogue code GREEN-044
\setcounter{section}{131}
\section{Collinear 4-tuples Force Collinear 5-tuples}\label{132}
{\small\sffamily UnsolvedMath ID 132 · Geometry}
\subsection{Definitions and mathematical statement}

\paragraph{Shared notation.}$\mathbb N=\{1,2,\ldots\}$, and $\mathbb Z,\mathbb Q,\mathbb R,\mathbb C$ denote the integers, rationals, reals and complex numbers. Any other conventions are specified locally.


For a finite set $A\subset\mathbb R^2$, define
\[
 t_4(A)=\bigl|\{B\subset A:|B|=4\text{ and all points of }B\text{ lie on one line}\}\bigr|.
\]
Thus $t_4$ counts unordered subsets of four distinct points. If no five points are collinear, each four-point subset corresponds to a line meeting $A$ in exactly four points; this matches the line-counting convention of the cited source.
\paragraph{Statement recorded in the supplied source.}
For every real $c>0$, is there an integer $n_0(c)$ such that every finite $A\subset\mathbb R^2$ with $|A|=n\geq n_0(c)$ and $t_4(A)\geq c n^2$ contains five distinct collinear points? Equivalently, if
\[
 T_4(n)=\max\{t_4(A):A\subset\mathbb R^2,\ |A|=n,\text{ no five points of }A\text{ are collinear}\},
\]
is $T_4(n)/n^2\longrightarrow0$ as $n\to\infty$? \sref{132}{2}
\subsection{Short English statement}
Must every large point set with a fixed positive quadratic number of collinear four-point subsets contain five collinear points?
\subsection{Sources}
\begin{description}
\item[\hypertarget{source:132:1}{S1}] ULAM.AI and the UnsolvedMath Contributors, \emph{UnsolvedMath: A Curated Collection of Open Mathematics Problems}, record 132 (GREEN-044). CC BY 4.0. \url{https://huggingface.co/datasets/ulamai/UnsolvedMath}.
\item[\hypertarget{source:132:2}{S2}] József Solymosi and Miloš Stojaković, \emph{Many collinear $k$-tuples with no $k+1$ collinear points}, Discrete \& Computational Geometry 50 (2013), 811--820, §1, Erdős conjecture for $r=5,k=4$. \url{https://arxiv.org/abs/1107.0327}.
\item[\hypertarget{source:132:3}{S3}] Ben Green, \emph{100 Open Problems}, maintained author list, Problem 71 and its comments, accessed 7 October 2026. \url{https://people.maths.ox.ac.uk/greenbj/papers/open-problems.pdf}.
\end{description}
\label{end:132}
\end{document}
